\section{Related Work}
\label{sec:related_work}

\inlinepar{Aggregation in the provenance semiring framework.}
Our approach builds on~\cite{amsterdamer2011provenance}, in which
\sql{HAVING} conditions become formal comparison expressions between
elements of a provenance semimodule; \cite{amsterdamer2011provenance}
stops at producing these expressions, whose evaluation, as well as
probability computation, is not considered. \cite{pintor2025dbms}
follows this approach with a DBMS-independent rewriting system that
stores the expressions as strings. The present work starts where these
two end, by giving the expressions a value in every commutative
m-semiring.

\inlinepar{Negation in the provenance semiring framework.}
Since provenance semirings \cite{green2007provenance,green2017provenance}
do not support non-monotone queries, \cite{geerts2010database} extended
them with a monus, yielding m-semirings. Some m-semirings lack axioms one
may expect~\cite{amsterdamer2011limitations}, notably the distributivity
of $\otimes$ over $\ominus$ discussed in Section~\ref{sec:preliminaries};
the model has nonetheless been fruitful, as the basis of a provenance for
SPARQL~\cite{DBLP:journals/jacm/GeertsUKFC16}, of a semiring
generalization of Codd's theorem~\cite{badia2025codd} and of the
\sql{EXCEPT} operator of ProvSQL~\cite{sen2026provsql}. Quotient
semirings of polynomials with dual
indeterminates~\cite{gradel2025provenance} are an alternative, less
amenable to implementation.

\inlinepar{Provenance systems.}
ProvSQL~\cite{sen2026provsql}, the system we extend, rewrites queries
over provenance-tracked relations and stores annotations as a circuit;
its aggregation follows~\cite{amsterdamer2011provenance}, recording
comparisons on aggregates as formal expressions that can be displayed but
not evaluated; it is also a probabilistic database.
GProM~\cite{arab2018gprom} is a rewriting middleware over several
backends that computes several forms of provenance without full semiring
semantics; it joins aggregation results with the provenance of each group
and treats \sql{HAVING} as a plain selection, so that it does not record
how the provenance depends on the comparison, which precludes its use as
a probabilistic database.

\inlinepar{Aggregation in probabilistic databases.}
\cite{re2009trichotomy} characterizes which aggregate conjunctive queries
with a single \sql{HAVING} atom can be evaluated or approximated in
polynomial time over probabilistic databases, by the shape of the query;
the base case of its safe plans, the convolution of the marginal
distributions of independent tuples in a small semiring, is what
Proposition~\ref{prop:agg-cmp-poly-prob} restates in our setting, whereas
Section~\ref{sec:algorithms} classifies combinations of semirings and
aggregates for symbolic provenance, which~\cite{re2009trichotomy} does
not consider. \cite{abiteboul2011capturing} studies
aggregation over probabilistic XML, for terminal aggregates only. Closest
to us, \cite{fink2012aggregation} extends the construction
of~\cite{amsterdamer2011provenance} for probability evaluation and
supports \sql{HAVING}, with a semantics specific to the Boolean case and
to aggregates over commutative monoids, implemented in a prototype of
SPROUT that was not released. Among the probabilistic database prototypes
MystiQ~\cite{DBLP:conf/sigmod/BoulosDMMRS05},
Trio~\cite{dblp:conf/vldb/agrawalbshnsw06}, Orion~\cite{singh2008orion},
and MayBMS~\cite{huang2009maybms} with its public SPROUT
evaluator~\cite{DBLP:conf/icde/OlteanuHK09}, none is maintained,
aggregation over uncertain data is at best partial
(Trio~\cite{murthy2011aggregation}, Orion), and none gives \sql{HAVING}
an uncertain semantics.
ProbLog~\cite{DBLP:journals/tplp/FierensBRSGTJR15,DBLP:conf/pkdd/DriesKMRBVR15},
a probabilistic programming system, supports aggregation and conditions
on aggregates through grounding and weighted model counting, but no other
(m-)semiring, and does not scale to databases of realistic
sizes~\cite{sen2026provsql}.

\inlinepar{Aggregation over compact representations.}
\cite{bakibayev2013aggregation} adds aggregation and \sql{HAVING} to a
factorized database engine, without provenance and in a deterministic
setting; whether factorized representations could compress the possible
worlds of our semantics is a natural question for future work.
